% ============================================================
%  Neural Network Foundations and Recurrent Neural Networks
%  Sources: Karpathy RNN Blog, Colah LSTM Blog,
%           Jurafsky & Martin Ch.9 (rnn_lstm.pdf),
%           Sequence Learning Slides (3_sequence_learning_pptx.pdf)
% ============================================================


% ============================================================
\section{Neural Network Foundations}
% ============================================================

\subsection{The Perceptron Formula}

A perceptron computes $f(wx + b)$, where $w$ is the weight, $x$ is the input, $b$ is the
bias, and $f$ is a non-linear activation function such as sigmoid, tanh, or ReLU. The
linear combination $wx + b$ is first computed, and then passed through $f$ to introduce
non-linearity.

% -----------------------------------------------------------
\subsection{A Neural Network as a Collection of Perceptrons}
% -----------------------------------------------------------

A neural network is essentially many perceptrons working together. Each node in a layer
is one perceptron. For example, a layer with 3 nodes contains 3 perceptrons, each
computing its own $f(wx + b)$.

% -----------------------------------------------------------
\subsection{How a Neural Network Learns Complex Non-Linear Patterns}
% -----------------------------------------------------------

A single perceptron applies one small non-linear transformation --- a single ``bend'' in
the data. This alone cannot capture complex patterns. However, when hundreds of these
bends are stacked across multiple layers, the combination becomes powerful.

Think of it like origami: one fold of paper is simple, but many folds in different
directions produce a complex 3D shape. Similarly:

\begin{itemize}
    \item Each neuron applies a tiny, simple transformation.
    \item A layer of neurons applies many transformations simultaneously.
    \item The next layer transforms those results again.
    \item Layer after layer, the network builds increasingly abstract representations.
\end{itemize}

For example, in image classification:
\begin{itemize}
    \item Layer 1 detects edges.
    \item Layer 2 combines edges into shapes.
    \item Layer 3 combines shapes into features.
    \item The final layer produces the class prediction.
\end{itemize}

No single step is intelligent. The depth and combination of simple transformations is
what enables the network to learn complex patterns.

% -----------------------------------------------------------
\subsection{How Backpropagation Works}
% -----------------------------------------------------------

Training a neural network follows four steps:

\begin{enumerate}
    \item \textbf{Forward pass:} Data passes through all layers and the network produces
          a prediction.
    \item \textbf{Loss calculation:} The prediction is compared to the true label. The
          loss is a single number representing how wrong the prediction was.
    \item \textbf{Backpropagation:} Starting from the loss, blame is traced backwards
          through the network using the \textbf{chain rule}. Each weight's contribution
          to the error is quantified as a gradient.
    \item \textbf{Gradient descent:} Each weight is nudged slightly in the direction that
          reduces the loss. The size of the nudge is controlled by the \textbf{learning
          rate}.
\end{enumerate}

This process repeats thousands of times. Over time, the weights converge to values that
produce accurate predictions.

% -----------------------------------------------------------
\subsection{Which Layers Get Their Weights Updated}
% -----------------------------------------------------------

In a network with 1 input layer, 2 hidden layers, and 1 output layer, weight updates
happen in all layers \textit{except} the input layer:

\begin{itemize}
    \item Output layer weights \checkmark
    \item Hidden layer 2 weights \checkmark
    \item Hidden layer 1 weights \checkmark
    \item Input layer weights \texttimes\ \ (the input layer has no weights --- it only
          passes data forward)
\end{itemize}

% -----------------------------------------------------------
\subsection{Backpropagation vs.\ Gradient Descent --- An Important Distinction}
% -----------------------------------------------------------

These two terms are often confused but refer to separate steps:

\begin{itemize}
    \item \textbf{Backpropagation:} Computes the gradients --- how much each weight
          contributed to the loss. It traces blame backwards using the chain rule.
    \item \textbf{Gradient descent:} Uses those gradients to actually update the weights
          --- nudging each weight to reduce the loss.
\end{itemize}

Backpropagation figures out \textit{how much} each weight is to blame. Gradient descent
does the actual correction.

% -----------------------------------------------------------
\subsection{When Backpropagation Begins}
% -----------------------------------------------------------

Backpropagation cannot start until a loss value exists. The sequence is:

\begin{enumerate}
    \item Forward pass completes and an output is produced.
    \item Loss is calculated by comparing the output to the true label.
    \item Backpropagation begins, working backwards from the loss.
\end{enumerate}

% -----------------------------------------------------------
\subsection{Input Layer Size for Tabular Data}
% -----------------------------------------------------------

For tabular data, the number of neurons in the input layer equals the number of features.
For example, a dataset with 3 columns (features) requires 3 neurons in the input layer
--- one per feature, all fed in simultaneously.

% -----------------------------------------------------------
\subsection{Stochastic Gradient Descent: Training Row by Row}
% -----------------------------------------------------------

For a tabular dataset with 10 rows, training with SGD proceeds as follows:

\begin{enumerate}
    \item Network processes row 1 (forward pass) $\to$ output produced $\to$ loss
          calculated $\to$ backpropagation $\to$ gradient descent (weights updated).
    \item Move to row 2 --- repeat the same process with the already-updated weights.
    \item Continue through row 10.
\end{enumerate}

This approach --- updating weights after every single sample --- is called
\textbf{Stochastic Gradient Descent (SGD)}.

Other approaches include:
\begin{itemize}
    \item \textbf{Batch Gradient Descent:} Process all rows first, then update weights
          once.
    \item \textbf{Mini-batch Gradient Descent:} Process a small group of rows, then
          update weights.
\end{itemize}

% -----------------------------------------------------------
\subsection{What Is an Epoch}
% -----------------------------------------------------------

One \textbf{epoch} is one complete pass through the entire training dataset. For a
10-row dataset, one epoch means the network has processed all 10 rows once.

